\chapter{Low Power: Deep Sleep and Retained State}
\label{ch:lowpower}\index{deep sleep}\index{low power}\index{RTC}

A board that runs from a battery spends most of its life asleep. The ESP32-S3's
deepest idle mode, \emph{deep sleep}, drops the whole chip to a handful of
microamps by powering down the digital core, the radios, and most of the SRAM. The
catch is what that costs: in deep sleep the core is not paused, it is \emph{off},
and waking is a chip reset. Your \texttt{Main} does not resume where it left off; it
starts again from the top. This chapter is about programming with that fact rather
than against it, through the \texttt{ESP32S3.RTC} driver and the
\texttt{esp32s3\_rtc\_sleep} example.

\section{What survives, and what does not}\index{retained memory}\index{RTC domain}

One island stays powered through deep sleep: the \emph{RTC domain}. It keeps the
real-time clock counting, holds a small block of \emph{slow RTC memory}, and can
hold a few pins in a fixed state. Everything else, the CPU registers, the main SRAM,
your task stacks, is gone on wake. The consequence is a clean rule:

\begin{quote}
Any state that must outlive a sleep lives in retained RTC memory. Everything else is
rebuilt from scratch on every wake.
\end{quote}

The driver exposes that memory as a small indexed array of 32-bit words:

\begin{lstlisting}
function  Read  (Index : Word_Index) return Interfaces.Unsigned_32;
procedure Write (Index : Word_Index; Value : Interfaces.Unsigned_32);
\end{lstlisting}

A boot counter, a state-machine phase, the next alarm time: these are what go here.
They survive the reset; a global variable does not.

\section{Why the board woke}\index{wake cause}

Because every wake re-runs \texttt{Main}, the first thing an application asks is
\emph{why} it is running: a fresh power-on, or a return from sleep? The driver
answers with the wake cause:

\begin{lstlisting}
type Wake_Cause is (Power_On, Deep_Sleep_Timer, Deep_Sleep_GPIO, Other_Reset);
function Last_Wake return Wake_Cause;
\end{lstlisting}

On \texttt{Power\_On} the application initialises its retained state from nothing. On
a \texttt{Deep\_Sleep\_*} wake it instead restores from retained memory and continues
the cycle. That branch, taken once at the top of \texttt{Main}, is the whole shape of
a low-power application.

\section{Going to sleep}\index{wake source}

Two wake sources cover most uses, a timer and a pin, and the driver gives one call
for each:

\begin{lstlisting}
procedure Deep_Sleep_For   (Wake_After : Duration);
procedure Deep_Sleep_Until (Pin : ESP32S3.GPIO.Pin_Id; High : Boolean := True);
\end{lstlisting}

\texttt{Deep\_Sleep\_For} arms the RTC timer and sleeps; the chip wakes after roughly
the requested duration, counted on the always-on RTC clock. \texttt{Deep\_Sleep\_Until}
arms an RTC GPIO and sleeps until that pin reaches the given level, which is how a
button or a sensor's interrupt line brings the board back. Neither call returns: the
next thing that runs is \texttt{Main}, from the top, after the wake reset.

The \texttt{esp32s3\_rtc\_sleep} example is the pattern in miniature. A boot counter
lives in retained memory. Each boot it reads the wake cause, bumps the counter, and
for the first few cycles sleeps on a timer. The counter climbing across the resets,
and the wake cause turning from \texttt{Power\_On} into \texttt{Deep\_Sleep\_Timer},
is the proof that retained memory, deep sleep, and timer wake all work together.
\texttt{esp32s3\_rtcio\_hold} shows the companion trick: holding an output pin at a
fixed level \emph{through} the sleep, so a load stays driven while the core is off.

\section{Holding pins, and the watchdog}\index{pin hold}\index{watchdog}

By default a GPIO floats when the core powers down, which is wrong for any pin that
must stay asserted, a chip-select held inactive, an enable line, a MOSFET gate. The
RTC can latch such a pin so it keeps its level across the sleep, and the application
releases the latch after the next wake once the driver has re-taken the pin.

One system timer can spoil all of this: the RTC super-watchdog, which resets the chip
if it is not fed. Deep sleep does not feed it, so the driver provides
\texttt{Disable\_Super\_Watchdog} to disarm it around a sleep. Forgetting this is a
classic cause of a board that resets a moment after it should have slept.

\section{Practical notes}

\cnote{ESP-IDF gives you \texttt{esp\_deep\_sleep\_start}, \texttt{esp\_sleep\_enable\_*}
wake-source calls, and the \texttt{RTC\_DATA\_ATTR} attribute to place a variable in
retained memory. The pieces here are the same pieces; they are just a small Ada
driver you can read end to end instead of a framework.}

Two things surprise people the first time. First, the \textbf{USB-Serial-JTAG console
drops on every sleep}, because that peripheral powers down with the core, so a board
that sleeps in a tight loop is hard to talk to; the \texttt{rtc\_sleep} example
deliberately stops sleeping after a few cycles and then repeats its final state so
the console can be captured. Second, a board left \textbf{cycling through sleep too
fast} can be genuinely hard to reflash, because it is awake for only a sliver of each
cycle; the recovery is to force the ROM download mode by hand (hold the BOOT button
while replugging) and flash a benign image. Sleep is a sharp tool, and the wake-reset
model (the same path the bootloader takes on any reset, Chapter~\ref{ch:bootloader})
is what makes it both cheap and unforgiving.
